\chapter{HGT 编码器与婚姻评分}\label{chap:hgt}

本章描述补全管线第一阶段的计算骨干：在第 \ref{chap:data} 章构建的异构图上训练异构图变换器（Heterogeneous Graph Transformer，HGT）\citep{hu2020hgt}，对每个候选对 \((m,w)\) 输出一个可比较的婚姻可信度分数。该分数既是后续 SEAL 子图模式抽取与多智能体协商的入口，也是消融实验中量化母系边贡献的统一指标。本章按"输入特征工程—HGT 卷积层—婚姻评分 MLP—训练循环—实验配置与定量结果"五节展开，所有数字均与 \texttt{src/config.py} 与 \texttt{src/model/hgt.py} 同源。

\section{节点输入特征工程}\label{sec:hgt-feature}

人物节点是评分任务唯一直接参与的节点类型。\texttt{PersonEmbedder} 把 DS0001 中按列分离的结构化字段重新组织为五条平行特征流，再统一投影到 \(d=128\) 维节点向量。表 \ref{tab:features} 列出五条流的来源、形状与处理方式。

\begin{table}[h]
\centering
\caption{\texttt{PersonEmbedder} 的五条平行输入流。}\label{tab:features}
\footnotesize
\begin{tabularx}{\linewidth}{@{}ll>{\raggedright\arraybackslash}X>{\raggedright\arraybackslash}X@{}}
\toprule
张量 & 形状 & 来源 & 处理方式 \\
\midrule
\texttt{x\_sex} & \([N]\) & DS0001 \texttt{SEX}（保留 1=女、2=男的原始编码） & 长度 3 词表的 \texttt{Embedding(3,8)}（索引 0 留给缺失） \\
\texttt{x\_relationship} & \([N]\) & DS0001 \texttt{RELATIONSHIP}，截断到出现频次最高的 64 类 & \texttt{Embedding(64,16)} \\
\texttt{x\_continuous} & \([N,1]\) & DS0001 \texttt{BIRTHYEAR} & 在非缺失行上估计均值方差后做列内 \(z\) 分数标准化，缺失行保持为 0 \\
\texttt{x\_occupational} & \([N,3]\) & DS0001 \texttt{POSITION}/\texttt{TITLE}/\texttt{SALARY} & 三个 0/1 指示位拼接 \\
\texttt{x\_macro} & \([K{=}5]\) & DS0009 粮价 + DS0011 灾荒 + 王朝纪元号 & 由队列年广播到 \(N\) 行 \\
\bottomrule
\end{tabularx}
\end{table}

五条流拼接后的总维度为 \(8+16+1+3+5=33\)。\texttt{PersonEmbedder} 借助一层线性投影将其映射到 \(\mathbb{R}^{128}\)，再依次作 GELU 激活与层归一化（Layer Normalization，LN）：
\begin{equation}\label{eq:embedder}
  h^{(0)}_v \;=\; \mathrm{LN}\!\Bigl(\mathrm{GELU}\!\bigl(W_{\text{person}}\,
    [\,e_{\text{sex}}\;\|\;e_{\text{rel}}\;\|\;x_{\text{cont}}\;\|\;x_{\text{occ}}\;\|\;x_{\text{macro}}\,]\bigr)\Bigr)
  \;\in\; \mathbb{R}^{128}.
\end{equation}
其中 \(W_{\text{person}}\in\mathbb{R}^{33\times 128}\)，\(e_{\text{sex}}\) 与 \(e_{\text{rel}}\) 分别由两张嵌入查找表给出，\(\|\) 表示按最后一维拼接。

\texttt{x\_macro} 的设计是该工程中关键的因果对齐细节。五项宏观协变量为 \(\textsf{cohort\_year\_z}\)（队列年的列向 \(z\) 分数）、\(\textsf{grain\_price\_z}\)（DS0009 粮价的列向 \(z\) 分数）、\(\textsf{grain\_price\_yoy}\)（粮价的逐年差分）、\(\textsf{era\_id}\)（清代纪元号：乾隆 0、嘉庆 1、道光 2、咸丰 3、同治 4、光绪 5、宣统 6）以及 \(\textsf{disaster\_flag}\)（DS0011 灾荒 0/1 指示）\citep{campbell2011data}。这五列以 \((n_{\text{years}}, K)\) 的查找表 \texttt{macro\_table} 附着在图对象上；当时间因果子图构造函数 \texttt{subgraph\_at\_year}\((t)\) 被调用时，它会读取第 \(t-1\) 行并将其挂到 \texttt{data["person"].x\_macro}。如此一来，无论快照年是哪一年，每次前向传播看到的宏观协变量都严格对齐到当前打分的队列年，避免把快照年的灾荒或粮价水平泄漏给更早的队列。具体而言，原本设计中曾使用的 \texttt{AGE\_IN\_SUI} 字段已被移除，因为它依赖于该人物最后一次观测年份，会把婚后状态泄漏给婚前打分；所需的年龄信息可由模型从 \(\textsf{cohort\_year\_z}\) 与 \texttt{BIRTHYEAR} 自行重建。

家户、社区、八旗三类节点采用相同的 \(\mathrm{Linear}\to\mathrm{GELU}\to\mathrm{LN}\) 投影到 128 维：家户节点输入 5 维（户口规模、夫妻数、世代数、户主任职指示、末次观测年的归一化值），社区节点输入 1 维（村庄人口的归一化），八旗节点输入与旗数等长的独热编码。这一设计使所有节点类型的隐表征位于同一线性空间，便于后续 HGT 卷积层做跨类型注意力。

\section{HGT 卷积层}\label{sec:hgt-encoder}

设节点 \(v\) 在第 \(\ell\) 层的隐向量为 \(h_v^{(\ell)}\)，节点类型函数为 \(\tau:V\to\mathcal{T}\)，关系类型集合为 \(\mathcal{R}\)。HGT 在每条边类型 \(r\in\mathcal{R}\) 上以多头注意力机制\citep{vaswani2017attention} 聚合邻居信息，并通过类型化的查询、键、消息投影区分不同节点类型与不同关系\citep{hu2020hgt}。注意力分数定义为
\begin{equation}\label{eq:hgt-att}
  \alpha_{u,v}^{r} \;=\; \mathrm{softmax}_u\!\left(
    \frac{\bigl(W^{q}_{\tau(v)}\,h_v^{(\ell)}\bigr)^{\!\top}\,W^{r}_{\mathrm{att}}\,\bigl(W^{k}_{\tau(u)}\,h_u^{(\ell)}\bigr)}
         {\sqrt{d/H}}
    \cdot \mu_{\tau(u),\,r,\,\tau(v)}\right),
\end{equation}
节点更新规则为
\begin{equation}\label{eq:hgt-upd}
  h_v^{(\ell+1)} \;=\; \mathrm{LN}\!\left(
    h_v^{(\ell)}\;+\;\mathrm{Dropout}_p\!\Bigl(
      \sum_{r\in\mathcal{R}}\sum_{u\in\mathcal{N}_r(v)}
        \alpha_{u,v}^{r}\,W^{r}_{\mathrm{msg}}\,W^{v}_{\tau(u)}\,h_u^{(\ell)}\Bigr)
    \right).
\end{equation}
上式中 \(W^{q}_{\tau(v)}\)、\(W^{k}_{\tau(u)}\) 与 \(W^{v}_{\tau(u)}\) 是按节点类型注册的查询、键、值线性投影；\(W^{r}_{\mathrm{att}}\) 与 \(W^{r}_{\mathrm{msg}}\) 是按关系注册的注意力与消息投影，用于将来源类型的隐空间转换到目标类型的隐空间；\(\mu \in \mathbb{R}^{|\mathcal{T}|\times|\mathcal{R}|\times|\mathcal{T}|}\) 是可学习的关系先验张量，给"特定来源类型—关系—目标类型"三元组提供一个全局乘性偏置。在本研究的图模式下 \(|\mathcal{T}|=4\)、\(|\mathcal{R}|=9\)，故 \(\mu\in\mathbb{R}^{4\times 9\times 4}\)。残差项 \(h_v^{(\ell)}\) 与外层 LayerNorm 共同稳定深层堆叠下的梯度，并在某种关系上没有任何邻居的节点处保留来自上一层的信号；Dropout 在训练阶段对聚合后的消息以概率 \(p\) 随机置零，缓解过拟合\citep{kipf2016gcn}。

超参取自 \texttt{src/config.py} 单一来源：隐维 \(d=128\)，深度 \(L=2\)，注意力头数 \(H=4\)，丢弃概率 \(p=0.2\)。基于 PyG\citep{fey2019pyg} 的 \texttt{HGTConv} 实现要求所有边类型在元数据快照中有显式注册，因而即便母系边在边消融实验中被置空，仍必须保留 \texttt{r\_ms} 与 \texttt{r\_md} 两个键、只把 \texttt{edge\_index} 替换为形状 \((2,0)\) 的张量；否则按关系参数化的 \(W^{r}_{\mathrm{att}}\)、\(W^{r}_{\mathrm{msg}}\) 会被裁剪，致使消融模型与未消融模型权重无法对齐。

\section{婚姻评分 MLP}\label{sec:hgt-scorer}

HGT 编码器输出每个人物节点的 128 维表征 \(h\) 后，由一个轻量的二元交互 MLP 给候选对 \((m,w)\) 打分。设拼接向量为 \([h_m\,\|\,h_w\,\|\,|h_m-h_w|\,\|\,h_m\odot h_w]\in\mathbb{R}^{4d}\)，其中 \(\odot\) 为哈达玛积，则评分函数定义为
\begin{equation}\label{eq:scorer}
  \psi(h_m,h_w) \;=\; W_2\,\mathrm{Dropout}_p\!\bigl(\mathrm{GELU}(W_1\,[\,h_m\,\|\,h_w\,\|\,|h_m-h_w|\,\|\,h_m\odot h_w\,])\bigr)\;\in\;\mathbb{R},
\end{equation}
其中 \(W_1\in\mathbb{R}^{128\times 512}\) 把 \(4d=512\) 维交互向量投到隐空间，\(W_2\in\mathbb{R}^{1\times 128}\) 输出标量 logit。\(4d\) 这一维度并非随意选取：原始两向量保留各自的全局位置，绝对差 \(|h_m-h_w|\) 编码两者在嵌入空间的对称差异，哈达玛积 \(h_m\odot h_w\) 编码逐维一致性；这四种二元交互在张量分解与对偶建模中已有成熟先例\citep{vaswani2017attention}，能在不引入外积矩阵的前提下保留线性无法捕捉的二阶相关。

\section{训练循环}\label{sec:hgt-train}

每个训练 epoch 先打散全部训练队列年的顺序，再以 \texttt{BATCH\_YEARS}\,=\,4 个队列年为粒度合成一次优化器步。对每个队列年 \(t\)，\texttt{subgraph\_at\_year}\((t,\;\text{drop}=\text{train}[t]\cup\text{all\_val}\cup\text{all\_test})\) 构造用于消息传递的因果子图。具体而言，\(\text{train}[t]\) 是当前打分的正例对，必须从消息传递图中扣除以避免标签泄漏；\texttt{all\_val} 与 \texttt{all\_test} 也被全局扣除——前者用于守住模型选择阶段的验证目标尚未被偷看，后者用于守住测试集在训练阶段从未对消息流动产生贡献。任何允许测试期婚姻边参与消息传递的设置都会让评估数字偏乐观，故采用全局扣除。

正例由当前队列年的真实婚姻对 \((m,w,1)\) 给出，负例按 \texttt{NEG\_PER\_POS}\,=\,4 的比例在同一队列年内通过替换女性候选生成 \((m,w',0)\)。这种"队列内负采样"的设计是数据驱动的：跨年负采样会让模型只需识别队列年差异即可大幅降低损失，使得 recall@1 等指标显著虚高，但对真实部署场景中"在同一年内的若干合理候选间排序"毫无信号增益。每个 mini-batch 上的损失为
\begin{equation}\label{eq:loss}
  \mathcal{L} \;=\; -\frac{1}{|\mathcal{B}|}\sum_{(m,w,y)\in\mathcal{B}}\Bigl[\,y\,\log\sigma(\psi(h_m,h_w))\,+\,(1-y)\,\log\!\bigl(1-\sigma(\psi(h_m,h_w))\bigr)\,\Bigr],
\end{equation}
其中 \(\sigma\) 为 sigmoid，\(\mathcal{B}\) 为该批次中所有正负样本三元组。表 \ref{tab:hyperparams} 汇总训练超参，全部值与 \texttt{src/config.py} 同步。

\begin{table}[h]
\centering
\caption{训练阶段超参（与 \texttt{src/config.py} 同源）。}\label{tab:hyperparams}
\footnotesize
\begin{tabularx}{\linewidth}{@{}l l >{\raggedright\arraybackslash}X@{}}
\toprule
名称 & 取值 & 说明 \\
\midrule
\texttt{HIDDEN} & 128 & 节点隐维 \\
\texttt{NUM\_LAYERS} & 2 & HGT 层数 \\
\texttt{HEADS} & 4 & 注意力头数 \\
\texttt{DROPOUT} & 0.2 & 节点更新与评分 MLP 共享 \\
\texttt{LR} & \(10^{-3}\) & AdamW 初始学习率 \\
\texttt{WEIGHT\_DECAY} & \(10^{-2}\) & 解耦权重衰减 \\
\texttt{BATCH\_YEARS} & 4 & 每步合并的队列年数 \\
\texttt{NEG\_PER\_POS} & 4 & 每正例配的队列内负样本数 \\
\texttt{GRAD\_CLIP} & 1.0 & 全局梯度范数裁剪上限 \\
\texttt{EPOCHS} & 30 & 训练总轮次 \\
\bottomrule
\end{tabularx}
\end{table}

优化器选用解耦权重衰减形式的 AdamW；学习率沿余弦曲线衰减至 0；前向反向均在原始浮点精度下进行，因为 PyG 的关系分块矩阵乘内核不接受混合精度输入。代码层面基于 PyTorch\citep{paszke2019pytorch} 与 PyG\citep{fey2019pyg} 实现。

\section{实验配置与定量结果}\label{sec:hgt-eval}

\textbf{数据切分。}训练—验证—测试切分由婚姻发生年驱动，而不是随机抽样。具体地，所有 \texttt{r\_hw} 边按其 \(\texttt{edge\_time}\) 字段所记录的婚姻年份分桶：年份不晚于 1855 的边进入训练集，1855 之后到 1879 之间的边进入验证集，1879 之后到 1909 之间的边进入测试集。这种基于队列年的切分确保任何时刻模型只能利用严格早于打分时间点的边，从而把"未来不可见"的因果约束嵌入数据本身。

\textbf{模型选择。}训练过程在每个 epoch 末尾都会在验证集上做一次按队列年的匈牙利匹配（Hungarian assignment）评估，得到当年所有男性候选与女性候选间的最优一对一指派，再统计 recall@1。该指标在历史档案补全任务中比 ROC-AUC 等指标更贴近实际诉求，因为补全的目标本身就是为每位男性找到最合适的妻子。验证 recall@1 最高的检查点被分别保存为 \texttt{checkpoints/best\_ablated.pt} 与 \texttt{checkpoints/best\_unablated.pt}。两个文件不可互换：消融图缺少 \texttt{r\_ms} 与 \texttt{r\_md} 两条关系，因此在消融条件下训练得到的 \(W^{r}_{\mathrm{att}}\) 与 \(W^{r}_{\mathrm{msg}}\) 在这两条关系上从未被更新；若直接载入未消融图则相当于在保留权重为随机初始化的关系上做推断，结果不再具有意义。

\textbf{定量结果。}本章节最终在 1880--1909 的留出测试集（共 131{,}140 个候选对）上完成评估，HGT 编码器与婚姻评分 MLP 共同给出 \(\mathrm{Hit@10}=0.624\) 与 \(\mathrm{MRR}=0.428\)。在母系边消融实验中，未消融上界相对于消融模型，在按队列年聚合的 Hungarian recall@1 上平均高出约 \(62.6\%\)。这一相对差距说明母系边在父系碎裂图上对婚姻边的可识别度仍有显著贡献，也为后续 SEAL 子图模式与多智能体协商所要补足的"母家信号"提供了量化基线。
